\par
\SourceLeaf{488}%
I have yet to add that I submitted a detailed memoir on this subject to His — — — a year ago, and that he, recognizing the importance of continuing the Danish surveys, at least immediately secured the possibility of doing so by ordering me to link Lüneburg to the southern points of that survey, some of which consisted of temporary signal points. This was done last autumn, so that the future continuation no longer depends on the survival of the points whose permanence is uncertain.

If Your — — approve of these ideas and consider them worth working to put into effect, I must assume that the intention is to carry out the undertaking as perfectly as possible and in a manner that can stand honorably alongside any other. With that in mind, I would gladly undertake the work if it were entrusted to me, and would regard an interruption of my purely astronomical work, perhaps lasting longer than one summer, as no sacrifice.

As for the expense required to extend the Danish degree survey through the Kingdom of Hanover, it is admittedly impossible to state it in advance with any accuracy, since it depends on various circumstances that cannot be foreseen. The weather, for instance, has a considerable influence on the duration and hence on the cost. I am unaware of the costs of any degree survey having been made public. Even if I could obtain information from the Danish astronomer about the costs of his survey so far, any estimate based on it for similar undertakings in the Kingdom of Hanover would be very uncertain because of the great differences in local conditions. In Hanover, for example, cartage for transporting people and instruments will account for a very substantial share of the total costs, while in Denmark, where cartage is supplied ex officio in natura, these costs almost entirely drop out of the reckoning. Still, all things considered, I believe that the sum of 1500 pounds sterling would suffice to cover all costs, and it goes without saying that an account would then be rendered.

In his degree survey, Professor \textsc{Schumacher} has, in addition to a couple of amanuenses, two officers of captain’s rank as assistants. Their work is to travel through the area beforehand, find suitable points for triangulation stations and, after making preliminary measurements there, propose them to Prof. S. for selection. They then see to the erection of signals wherever necessary at the chosen points and to any other arrangements needed. In short, they make all the preparations so that the astronomer can proceed to the observations at every station without losing much time. Finally, where necessary, they take charge of the auxiliary operations during the observations. The qualities needed in such assistants are therefore less a particularly deep understanding of mathematics or astronomy than a lively zeal for the work, the utmost punctuality, a concern for the highest accuracy, a certain practical aptitude, some knowledge of construction, and some familiarity with the procedures of the authorities in our country with whom one deals in such matters. The Hanoverian survey would likewise need no more staff than the Danish one, and as soon as a decision is made I would make it my immediate concern to find suitable assistants myself. It would naturally be for Your — — to decide whether, \textit{in the respect mentioned above}, it might be advisable to let several more people attend the measurements as volunteers in order to prepare themselves for more exacting geodetic work. — —

Göttingen, May 30, 1819.
\begin{flushright}\textsc{Carl Friedrich Gauss}.\end{flushright}
